\section{Applying the framework to a central bank release programme}\label{sec:apply}

This section sets out how a central bank can apply the framework to its own
releases. It assumes no reading of the derivations. Each step names what the
institution must decide, who decides it, what evidence the step produces, and
where in the paper the method is justified. Appendix \ref{app:worksheets} holds
blank worksheets for each step, and Table \ref{tab:record} shows them filled for
the synthetic credit register.

\subsection{Roles}\label{sec:apply-roles}

Four roles are needed, and the first and third should not be held by the same
person.
\begin{description}
\item[Release owner.] A named senior officer, for example the chair of a
statistics release committee, who sets the tolerance, signs the decision and owns
the residual risk.
\item[Data owner.] The business area that collects the source register and knows
which products are drawn from it, how often, and under what legal basis.
\item[Disclosure analyst.] The statistician who declares the scenarios, runs the
attack simulation and computes the bounds. The analyst owns the correctness of the
calculation, not the decision.
\item[Reviewer.] A person outside the production line, in the risk, audit or
methodology function, who checks that the scenarios are plausible and that the
falsification tests of Section \ref{sec:falsify} were run.
\end{description}
The analyst and the data owner sit in the first of the three lines of defence,
the reviewer in the second, and internal audit, which checks the filed record
and the audit controls of Section \ref{sec:security}, in the third (Section
\ref{sec:gov}).


\subsection{Seven steps}\label{sec:apply-steps}

Figure \ref{fig:protocol} shows the seven steps as a decision process, and
Algorithm \ref{alg:protocol} states them compactly. Two loops matter. The dotted
loop is the re-assessment at each new vintage or at the expiry of the decision,
which starts again at the planning step and draws on what remains of the
programme's error budget. The dashed loop is a change of product: when no setting
meets the tolerance, as for loan-level files against an adversary who knows loan
histories (Section \ref{sec:eval-programme}), the institution moves to another
product, such as aggregate tables or access in a secure setting, and assesses it
from the start.

\begin{figure}[t]
\centering
\begin{tikzpicture}[font=\small,node distance=0.45cm,
  stp/.style={draw,rounded corners=2pt,align=center,minimum width=7cm,minimum height=0.7cm},
  gate/.style={draw,diamond,aspect=2.2,align=center,inner sep=1pt,fill=gray!15},
  >=Stealth]
  \node[stp] (s1) {1. List the programme: releases, horizon,\\access route, legal basis, protected parties};
  \node[stp,below=of s1] (s2) {2. Declare the scenarios, before any audit};
  \node[stp,below=of s2] (s3) {3. Set the tolerance, owner, $\alpha$ and expiry};
  \node[stp,below=of s3] (s4) {4. Plan the measurement: order of settings,\\sample sizes, error shares};
  \node[stp,below=of s4] (s5) {5. Measure per release and pooled,\\on fresh targets};
  \node[stp,below=of s5] (s6) {6. Test the model behind the number};
  \node[gate,below=0.6cm of s6] (g) {every declared\\bound meets\\its tolerance?};
  \node[stp,minimum width=3.4cm] (rel) at ($(g.south)+(-2.4,-1.2)$) {7. Release and file\\the decision record};
  \node[stp,minimum width=3.4cm,pattern=north east lines,pattern color=gray!40] (no) at ($(g.south)+(2.4,-1.2)$) {Do not release;\\change the product};
  \draw[->] (s1)--(s2); \draw[->] (s2)--(s3); \draw[->] (s3)--(s4); \draw[->] (s4)--(s5); \draw[->] (s5)--(s6); \draw[->] (s6)--(g);
  \draw[->] (g.west) -| node[near start,above]{yes} (rel.north);
  \draw[->] (g.east) -| node[near start,above]{no} (no.north);
  \draw[->,dashed] (no.east) -- ++(0.5,0) |- node[near start,right,align=left]{\scriptsize new product,\\\scriptsize new audit} (s1.east);
  \draw[->,densely dotted] (rel.west) -- ++(-0.5,0) |- node[near start,left,align=right]{\scriptsize new vintage\\\scriptsize or expiry} (s4.west);
\end{tikzpicture}
\caption{The assessment of a release programme as a decision process. Solid
arrows follow the seven steps of Section~\ref{sec:apply} and Algorithm~\ref{alg:protocol}; the dotted arrow is the
re-assessment at each new vintage or at expiry, which draws on the programme's
remaining error budget; the dashed arrow is a change of product, such as moving
from loan-level files to aggregates or to secure access, which needs a new audit.}
\label{fig:protocol}
\end{figure}

\begin{algorithm}[t]
\caption{Assessing a release programme}\label{alg:protocol}
\begin{algorithmic}[1]
\Require a release programme and the institution's governance record $G$
\State list the releases, horizon, access route, legal basis, protected parties \Comment{Step 1}
\State declare every scenario, pooled ones included, before any audit \Comment{Step 2}
\State set $\tau$ with its owner, harm class, $\alpha$ and expiry \Comment{Step 3}
\State fix the order of settings, the sample sizes and the error shares \Comment{Step 4}
\State measure per release and pooled on fresh targets; choose the setting by
  Algorithm~\ref{alg:selection} or Algorithm~\ref{alg:chain} \Comment{Step 5}
\State run the falsification tests; if one fails, stop: the figure is invalid \Comment{Step 6}
\If{every declared bound meets its tolerance}
  \State release, and file the decision record \Comment{Step 7}
\Else
  \State do not release; change the product and return to Step 1
\EndIf
\State at each new vintage or at expiry, re-assess from Step 4 on the remaining budget
\end{algorithmic}
\end{algorithm}



\paragraph{Step 1. List the programme, not the release.}
Start from the source register, not from the product about to be published. List
every product drawn from it over the planning horizon: monthly or quarterly
aggregates, supervisory publications, research extracts, data shared with other
authorities, and answers to ad hoc requests. For each product record its
frequency, the population it covers, the variables it publishes, and whether any
identifier or stable ordering persists from one vintage to the next. A persistent
pseudonym, a stable row order or a unique combination of rarely changing
attributes lets an adversary match records across vintages, which is the condition that Proposition \ref{prop:nopool} requires. With a persistent pseudonym the products are, for the adversary, one release on the union of their keys (Remark \ref{rem:finer}). Declare a horizon $m$,
the number of vintages the assessment covers. Worksheet 1 records the inventory.

\paragraph{Step 2. Declare the scenarios.}
For each product, and for each combination of products an outsider could hold
together, write down who the adversary is, what it knows, which records it
targets and what harm follows if it succeeds (Section \ref{sec:experiment}). For a
pooled scenario, say whether the adversary knows when the target's record is
present in each vintage, and why (Section \ref{sec:eval-threat}). The
declaration should be specific enough that someone can disagree with it. Table
\ref{tab:apply-map} gives starting declarations for release types common in
central banks. They are illustrations to be adapted, not a statement of what any
institution does. Worksheet 2 records each scenario.

\paragraph{Step 3. Set the tolerance.}
For each harm class the release owner sets a tolerance $\tau$, an error level
$\alpha$, an expiry date and the horizon $m$ of Step 1, and records the legal
basis, the parties the tolerance protects and the access route of the product
(Definition \ref{def:thresh}). The tolerance is the institution's written reading
of when identification is reasonably likely for that adversary and route. Two routes lead to a number. The direct route states the
largest probability of successful exposure the institution will accept for that
harm, taking account of its confidentiality duty and the public value of the
release. The indirect route states the relative weight the institution places on
harm against lost utility and reads the implied tolerance off the risk--utility
frontier (Proposition \ref{prop:tau}, Table \ref{tab:tau}). Either way, the
tolerance is written down before the risk is measured, and a programme tolerance
is set alongside the per-release one. Where the institution has not yet agreed a
tolerance for a harm, the record says so and no decision is made on that harm,
as in the last rows of Table \ref{tab:record}.

\paragraph{Step 4. Plan the measurement.}
Decide how many simulated attacks each scenario needs before running any. With no
observed success, certifying a tolerance $\tau$ at error level $\alpha$ needs
about $\log(1/\alpha)/\tau$ independent attacks, roughly $3/\tau$ at
$\alpha = 0.05$ \citep{hanley1983}; Table \ref{tab:cert} gives exact counts. A
tolerance of $10^{-3}$ needs 2{,}995 attacks, and one of $10^{-5}$ needs
299{,}572, so a very small tolerance can be expensive to certify by simulation.
If the release must pass every one of a fixed set of scenarios or subgroups, each
may be assessed at level $\alpha$, because that is an intersection--union test
(Section \ref{sec:gov}); divide $\alpha$ among them only when their bounds are to
be reported as one simultaneous statement. If a level such as a minimum cell
count is to be chosen from a grid, fix the order in which the levels are tested
before any result is seen (Proposition \ref{prop:fixedseq}). If scenarios may be added after results
are seen, for example because a reviewer proposes a stronger adversary, use a
registered spending sequence rather than a fixed split (Section
\ref{sec:sequential} and the companion article).

\paragraph{Step 5. Measure, per release and pooled.}
Run the declared attacks on the release as it would be published, on fresh
targets, and compute the exact upper bound $\overline{R}_{1-\alpha}$ for each
scenario (Section \ref{sec:est}). Then run the pooled attack over the programme
of Step 1, holding the same auxiliary information, and compute its bound in the
same way (Section \ref{sec:comp}). Report separately any pooled attack that adds
knowledge, such as more readings or the target's history across vintages. The
attacks must be reported side by side. Tables \ref{tab:decomp} and
\ref{tab:robust} show why: on the synthetic register, pooling with one reading
adds nothing unless the adversary knows the loan's history, and part of the
remaining excess comes from the adversary holding more readings, and a committee
needs to know which it is protecting against. Where a level is chosen from a
grid, audit the chosen level again on fresh targets. Report what the chosen
level keeps in volume, such as the share of exposure and of large borrowers, as
well as the share of records.

\paragraph{Step 6. Test the model behind the number.}
Before the bound is used, run the falsification tests of Section
\ref{sec:falsify}: calibration of the linkage posterior, the population-frequency
model, a stronger adversary, monotonicity of the mechanism, the composition test
and temporal transport. A failed test makes the assessment invalid, not failed:
the institution cannot release on it and must repair the model first.

\paragraph{Step 7. Decide, record and re-assess.}
Apply rule \eqref{eq:admiss}: the release passes a scenario only if the upper
bound, not the point estimate, is at or below the tolerance, and it passes
overall only if it passes every declared scenario, pooled ones included. Where
the programme fails, the release fails, even if it passes on its own. File the
decision record (Worksheet 3). Re-assess when the programme is about to pass its
horizon, when a new product is drawn from the same register, when a new public
source could serve as auxiliary information, when the mechanism changes, and at
expiry.

\subsection{Release types common in central banks}\label{sec:apply-map}

Table \ref{tab:apply-map} maps six release types to the objects of Steps 2 to 5.
The first column of the table is the product; the others say whom the adversary
targets, what it plausibly holds, which harm is at stake, which measure of this
paper applies, and where a pooled scenario arises.

\begin{table}[!t]
\centering
\caption{Starting declarations for release types common in central banks. The
entries are illustrations to be adapted to each institution's products, legal
basis and users; they do not describe the practice of any institution.}
\label{tab:apply-map}
\footnotesize
\setlength{\tabcolsep}{3pt}
\begin{tabular}{@{}>{\raggedright\arraybackslash}p{2.3cm}>{\raggedright\arraybackslash}p{2.4cm}>{\raggedright\arraybackslash}p{2.6cm}>{\raggedright\arraybackslash}p{2.3cm}>{\raggedright\arraybackslash}p{2.6cm}>{\raggedright\arraybackslash}p{2.6cm}@{}}
\toprule
Release type & Targets & Adversary holds & Harm & Measure in this paper & Pooled scenario \\
\midrule
Aggregates from a credit register (by sector, region, size) & Borrowers in small cells; lenders in concentrated cells & A firm's sector, region and size from public filings; a lender's own book & Identification of a borrower; a lender's position revealed & Prosecutor risk by cell (Section \ref{sec:three}); residual disclosure from totals (Appendix \ref{app:derivations}) & Monthly vintages; differencing of consecutive months \\
Loan-level research extract & Every borrower in the extract & Key fields and a noisy reading of exposure from other sources & Identification; attribute disclosure of arrears & Linkage attack of Section \ref{sec:eval-threat}; attribute rule & Successive extracts for the same project; extracts for different projects \\
Supervisory statistics by bank group & Individual banks & Each bank's own figures; published annual reports & A bank's confidential figure recovered & Contribution from totals; dominance and $p\%$ rules \citep{hundepool2012} & Quarterly series; tables that share a margin \\
Payment and transaction statistics & Merchants or payers in thin cells & Merchant category and location; own transactions & Identification; commercial information revealed & Prosecutor risk by cell; minimum contributor count & Daily or weekly series of the same cells \\
Household or firm survey microdata & Respondents & Demographic or firm characteristics from registers or social media & Identification; attribute disclosure of wealth or debt & Journalist and marketer measures (Section \ref{sec:three}); frequency model (Section \ref{sec:est}) & Survey waves; the public file and a scientific-use file \\
Answers to ad hoc data requests & Units in the requested breakdown & Everything already published from the same register & Identification by differencing & Differencing posterior (Appendix \ref{app:derivations}) & Each answer pooled with every earlier publication \\
\bottomrule
\end{tabular}
\end{table}

\subsection{When a release fails}\label{sec:apply-fail}

A failed scenario does not end the matter. It tells the committee what to change,
and the record should state which remedy was chosen and why. The remedies fall
into five groups.
\begin{enumerate}
\item \emph{Change the key.} Coarsen the published breakdown, or withdraw the
field that the adversary uses most (Table \ref{tab:ablation} shows how far each
field moves the risk on the synthetic register).
\item \emph{Raise the floor.} Increase the minimum cell count or add a dominance
rule, choosing the level by fixed-sequence testing on the upper bound of the
pooled scenario and confirming it on fresh targets (Table \ref{tab:robust}).
\item \emph{Break the link between vintages.} Rotate pseudonyms, change the row order,
or publish some products less often, so that the pooled adversary cannot match
records across vintages. Rotation is a control on linkability, not
pseudonymisation in the legal sense, and the pseudonym should be a keyed hash
under a key held apart from the release team. It helps only if the published
fields cannot relink the records; on the synthetic credit register the key and
the exact exposure relink every loan, and with exposure rounded to one
significant figure about one loan in ten (Section \ref{sec:eval-programme}).
\item \emph{Add noise.} Use a formal mechanism and account for its budget over the
horizon, remembering that its guarantee bounds posterior odds rather than naming a
linkage probability (Section \ref{sec:dp}).
\item \emph{Change the access route or the product.} Move the product from public
or scientific-use release to a secure research environment, where the risk is
governed by access controls and output checking rather than by the release alone
\citep{desai2016}, or publish aggregates in place of loan-level records. On the
synthetic credit register this is the realistic remedy against an adversary who
knows loan histories, because the floors that meet the tolerance keep a small
share of the exposure (Section \ref{sec:eval-programme}).
\end{enumerate}
Shortening the horizon $m$ is also a remedy, but only if the institution is prepared
to stop the series or re-assess it at the new horizon.

\subsection{Starting small}\label{sec:apply-pilot}

An institution does not need to apply the framework to every product at once. A
workable start is one register and one product family over one quarter:
\begin{enumerate}
\item choose the register whose products are published most often;
\item complete Worksheet 1 for it and agree a horizon;
\item agree one tolerance, for identification, with a named owner;
\item run one per-release and one pooled attack, with the sample sizes of Step 4;
\item file one decision record and present it to the release committee beside the
current method, so the committee can see where the two disagree.
\end{enumerate}
The pipeline that produces this paper runs the attacks and
writes the record. It needs, for each product, the published file or its
specification, the key fields, the auxiliary variables the adversary is declared to
hold and the tolerance. Section \ref{sec:real} shows the demographic path of the
code applied to two public files, which is the pattern for adapting it to an
institution's own data; the credit path has so far run only on the synthetic
register.

\subsection{What the framework does not decide}\label{sec:apply-limits}

The framework makes a decision reproducible; it does not make it for the
institution. It cannot say whether a declared adversary is realistic, whether a
tolerance is legally or ethically right, or whether a release is worth its risk.
Those remain judgements for the release owner. What it adds is that each judgement
is written down, owned and dated, and that the number it is compared with carries
its uncertainty and covers the whole programme.
